\documentclass[11pt]{article}

\usepackage{series}
\usepackage{amsmath,amssymb,amsthm}
\usepackage{booktabs}
\usepackage{hyperref}

\newcommand{\Tr}{\operatorname{Tr}}
\newcommand{\TV}{d_{\mathrm{TV}}}
\newcommand{\osc}{\operatorname{osc}}

\title{Born-Compatible Record Measures and the Classical Concentration Limit:\\
\large Separate Theorems and Their MTT Interface}

\author{Peter Nero}

\date{July 2026, Version 2}

\begin{document}
\maketitle

\begin{abstract}
The Born rule and the emergence of classical behavior are related, but they are
not the same mathematical problem. The first fixes outcome probabilities for a
declared quantum preparation and instrument. The second asks when retained
records and observables are well approximated by one stable classical
alternative. We formulate their interface without identifying them. Standard
Gleason-type results recover trace probabilities only after additive or
effect-noncontextual probability assumptions are supplied; they do not select
an MTT source measure. The current q79 program now provides a stronger,
domain-specific result: its canonical binary one-anchor Fock recorder emits an
exact stopped-output measure and second-moment capture descent from the selected
normal state, without an added Born axiom, stochastic primitive, fit, or
observed probability on that domain. General apparatus contexts and objective
single-history actualization remain open. Separately, we prove exact
concentration and persistence bounds. If one record has probability at least
\(1-\varepsilon\), the law is within \(\varepsilon\) in total variation of a
deterministic record, bounded observables differ by at most
\(\varepsilon\) times their oscillation, and a record with per-step escape
probability at most \(\eta\) survives \(n\) steps with probability at least
\(1-\varepsilon-n\eta\). These statements define a controlled classical limit.
They do not derive the outcome weights or select one realized history.
\end{abstract}

\section*{Version 2 Revision Note}
\addcontentsline{toc}{section}{Version 2 Revision Note}
\begin{description}
\item[Supersedes] \emph{Why the Born Rule and the Classical Limit Are the Same
Problem: A Projection-Based Shadow Bridge in Modal Triplet Theory}, version 1.
\item[Reason] The earlier paper inferred squared-norm basin weights from
projection, identified probability generation with classical concentration,
and described both as a unified resolution. It did not state the measure and
noncontextuality assumptions behind Gleason-style trace representation, and it
did not distinguish record stability from selection of one outcome.
\item[Resolution] This version separates the quantum probability source,
ordinary instrument dynamics, decoherence, concentration, record persistence,
and one-history actualization. It imports Gleason-type representation at its
proper conditional tier, records the exact current q79 binary-recorder result,
and proves finite concentration and persistence bounds independently.
\item[Retained result] Born-compatible record weights and the classical limit
can use the same outcome algebra and basin labels. Once a valid record law is
known, concentration into one robust record gives a controlled deterministic
approximation.
\item[Remaining boundary] MTT has not yet derived the required outcome law for
every allowed apparatus context, a universal concentration regime, or an
objective rule selecting one ontic history. Projection alone supplies none of
these three results.
\end{description}

\tableofcontents

\section{The corrected relation}

There are at least four distinct questions in a measurement process.

\begin{enumerate}
\item Which alternatives can become physical records?
\item What probabilities are assigned to those records?
\item When are interference terms between records operationally negligible?
\item When is one record so dominant and stable that a deterministic classical
description is accurate?
\end{enumerate}

A fifth question may be asked: why is one individual record actual rather than
another? Standard operational quantum mechanics need not answer that question
in order to predict recorded frequencies. Any ontic completion must answer it
without confusing the answer with the probability law.

The old title compressed questions 2 and 4 into one problem. The corrected
claim is narrower and more useful:
\[
\boxed{
\begin{gathered}
\text{Born-compatible record law}
+\text{ concentration and stability}\\
\Longrightarrow
\text{ controlled classical predictions}
\end{gathered}
}
\]
Each input has its own hypotheses and can fail independently.

\section{Measurement as an ordinary physical instrument}

Let \(\mathcal{H}\) be a complex Hilbert space and \(\rho\) a density operator.
A finite quantum instrument is a family
\[
\{\mathcal{I}_i\}_{i\in I}
\]
of completely positive, trace-nonincreasing maps such that
\(\sum_i\mathcal{I}_i\) is trace preserving. Its effects are
\[
E_i=\mathcal{I}_i^*(\mathbf{1}),\qquad
E_i\geq0,\qquad
\sum_iE_i=\mathbf{1}.
\]
The operational outcome law and conditional post-measurement state are
\begin{equation}
p_i=\Tr(\rho E_i),\qquad
\rho_i'=\frac{\mathcal{I}_i(\rho)}{p_i}
\quad(p_i>0).
\label{eq:instrument}
\end{equation}

Nothing in \eqref{eq:instrument} makes measurement metaphysically privileged.
An apparatus is a physical interaction that amplifies alternatives into
records. The instrument formalism records the input-output statistics and state
update of that interaction.

Suppose an upper model has measurable record regions \(B_i\) and an upper
probability law \(\mu_{\rho,\mathcal{I}}\) for the preparation and apparatus
context. The exact compatibility equation is
\begin{equation}
\mu_{\rho,\mathcal{I}}(B_i)=\Tr(\rho E_i)
\qquad\text{for every }i.
\label{eq:basin-trace}
\end{equation}
Writing down normalized basin weights does not prove
\eqref{eq:basin-trace}. The source law, the record regions, and the equality to
the trace weights all require independent construction.

\section{What Gleason-type theorems establish}

Gleason's theorem starts with a probability measure on the closed subspaces, or
equivalently projections, of a real or complex Hilbert space of dimension
greater than two. Countable additivity on mutually orthogonal subspaces implies
that the measure has the trace form
\[
\mu(P)=\Tr(\rho P)
\]
for a positive trace-class operator \(\rho\) \cite{Gleason1957}. Extensions
using positive-operator-valued measurements recover the trace form for
generalized effects and can include two-dimensional systems
\cite{CavesEtAl2004}.

These are representation theorems. Their assumptions already include a
probability assignment satisfying strong consistency conditions. They show the
form that such an assignment must take; they do not derive an upper physical
measure from non-injective projection.

For an MTT application, the logical order is therefore:
\[
\begin{gathered}
\text{selected preparation and record measure}\\
+\text{additivity or effect noncontextuality}
\end{gathered}
\Longrightarrow
\text{trace representation},
\]
followed by a separate proof of the basin--trace equality
\eqref{eq:basin-trace}. This separation prevents a conditional reconstruction
from being reported as a source theorem.

\section{Current q79 Born status}

\subsection{The exact canonical domain}

The current q79 program closes a specific operational domain. For the canonical
binary one-anchor recorder:
\begin{itemize}
\item the finite Hilbert/state/observable data and reduced dynamics are fixed;
\item the commuting nondemolition Fock output algebra supplies the record
events;
\item the selected normal state emits the stopped output measure;
\item second-moment capture descent is exact; and
\item no separate Born axiom, stochastic primitive, observed probability, or
numerical fit is added on that domain.
\end{itemize}

Thus equation \eqref{eq:basin-trace} has a selected operational realization for
that binary apparatus context. This is stronger than the conditional
Gleason-only status of the earlier paper.

\subsection{The quantifier boundary}

The canonical result does not yet prove:
\begin{itemize}
\item the same descent for every allowed apparatus and preparation context;
\item controlled corrections for finite-bandwidth or non-Markov detectors;
\item a pre-quantum probability semantics, if one is demanded;
\item or objective selection of one ontic history.
\end{itemize}

The correct status is therefore tiered:
\[
\begin{array}{ll}
\text{canonical q79 binary one-anchor output law:}&\text{exact},\\
\text{universal physical apparatus family:}&\text{open},\\
\text{objective single-history actualization:}&\text{open}.
\end{array}
\]

\section{Decoherence is not concentration}

Let \(\mathcal{D}\) be dephasing in a pointer decomposition. If
\[
\frac12\|\rho-\mathcal{D}(\rho)\|_1\leq\delta,
\]
then every effect \(0\leq E\leq\mathbf{1}\) satisfies
\begin{equation}
\left|
\Tr(\rho E)-\Tr(\mathcal{D}(\rho)E)
\right|
\leq\delta.
\label{eq:decoherence-bound}
\end{equation}
This is the operational meaning of approximate decoherence for the declared
effect family: coherences alter probabilities by at most \(\delta\).

Equation \eqref{eq:decoherence-bound} does not imply that one diagonal weight is
near one. The state
\[
\frac12|0\rangle\langle0|
+\frac12|1\rangle\langle1|
\]
is exactly decohered and maximally nonconcentrated on its two pointer records.
Decoherence helps explain stable alternatives and suppression of interference
\cite{Zurek2003}; a separate concentration estimate is needed for an
approximately deterministic record.

\section{The finite classical concentration theorem}

Let \(I\) be a finite record set and let
\[
p=(p_i)_{i\in I}
\]
be a probability law. For a function \(f:I\to\mathbb{R}\), define
\[
\osc(f)=\max_{i\in I}f(i)-\min_{i\in I}f(i).
\]

\begin{theorem}[Classical concentration]
Suppose a record \(i_\star\) satisfies
\[
p_{i_\star}\geq1-\varepsilon
\qquad(0\leq\varepsilon\leq1).
\]
Then
\begin{align}
\TV(p,\delta_{i_\star})
&=1-p_{i_\star}\leq\varepsilon,
\label{eq:tv}\\
\left|
\sum_{i\in I}p_i f(i)-f(i_\star)
\right|
&\leq\varepsilon\,\osc(f)
\label{eq:observable}
\end{align}
for every real function \(f\) on \(I\).
\end{theorem}

\begin{proof}
Using the convention
\(\TV(p,q)=\frac12\sum_i|p_i-q_i|\),
\[
\TV(p,\delta_{i_\star})
=\frac12\left(
1-p_{i_\star}+\sum_{i\neq i_\star}p_i
\right)
=1-p_{i_\star}.
\]
Also,
\[
\sum_i p_i f(i)-f(i_\star)
=\sum_{i\neq i_\star}
p_i\bigl(f(i)-f(i_\star)\bigr).
\]
Taking absolute values and using
\(|f(i)-f(i_\star)|\leq\osc(f)\) gives
\eqref{eq:observable}.
\end{proof}

This theorem gives an exact error budget for replacing the record law by the
deterministic prediction \(i_\star\). It does not say why
\(p_{i_\star}\) is large. That must follow from the preparation, dynamics,
environment, control regime, or an MTT source theorem.

\section{Record persistence}

Concentration at one time is not enough. A classical record should persist.
Let \(X_0,X_1,\ldots\) be a discrete-time record process on \(I\). No Markov
assumption is needed for the following conditional escape bound.

\begin{theorem}[Finite-horizon persistence]
Assume
\[
\Pr(X_0=i_\star)\geq1-\varepsilon
\]
and, for every \(k<n\),
\[
\Pr\!\left(
X_{k+1}\neq i_\star
\mid X_0=\cdots=X_k=i_\star
\right)
\leq\eta.
\]
Then
\begin{equation}
\Pr(X_0=\cdots=X_n=i_\star)
\geq(1-\varepsilon)(1-\eta)^n
\geq1-\varepsilon-n\eta.
\label{eq:persistence}
\end{equation}
\end{theorem}

\begin{proof}
The chain rule for conditional probabilities gives the first inequality.
Bernoulli's inequality gives \((1-\eta)^n\geq1-n\eta\). Finally,
\[
(1-\varepsilon)(1-n\eta)
\geq1-\varepsilon-n\eta.
\]
\end{proof}

Equation \eqref{eq:persistence} separates two classicality controls:
\[
\varepsilon
=\text{initial nonconcentration},
\qquad
\eta
=\text{per-step record escape}.
\]
A useful classical regime requires both to be small over the physical
observation horizon.

\section{How the two theorems fit together}

\begin{center}
\begin{tabular}{p{0.27\textwidth}p{0.31\textwidth}p{0.31\textwidth}}
\toprule
Layer & Required object & What it establishes\\
\midrule
Alternatives
& instrument effects and record algebra
& available records\\
Probability source
& selected state/measure and capture map
& the values \(p_i\)\\
Trace representation
& additivity or effect noncontextuality
& \(p_i=\Tr(\rho E_i)\)\\
Decoherence
& suppression of off-diagonal influence
& classical mixture approximation\\
Concentration
& \(p_{i_\star}\geq1-\varepsilon\)
& deterministic prediction with error \(\varepsilon\)\\
Persistence
& escape control \(\eta\)
& stable record over a finite horizon\\
Actualization
& interpretation or additional dynamics
& why one history is individually actual\\
\bottomrule
\end{tabular}
\end{center}

The same record labels and effects can appear in every row. That shared
interface is the legitimate connection between the Born law and the classical
limit. The logical arrows are not reversible. Concentration cannot derive the
general Born weights, and a Born law need not be concentrated.

\section{A precise MTT completion target}

To extend the canonical q79 result into a general Born-and-classical theorem,
MTT must supply a family indexed by physical contexts \(c\) and a classical
regime parameter \(L\):
\[
\bigl(
\rho_{c,L},
\{\mathcal{I}_{i,c,L}\},
\mu_{c,L},
\{B_{i,c,L}\}
\bigr).
\]
The required statements are:

\begin{enumerate}
\item \textbf{Source equality.}
\[
\mu_{c,L}(B_{i,c,L})
=\Tr(\rho_{c,L}E_{i,c,L})
\]
for every allowed context and outcome, with stated errors for nonideal
detectors.
\item \textbf{Concentration regime.} There is a selected
\(i_\star(c,L)\) and an explicit \(\varepsilon(c,L)\to0\) such that
\[
\mu_{c,L}(B_{i_\star,c,L})\geq1-\varepsilon(c,L).
\]
\item \textbf{Persistence regime.} A derived escape bound
\(\eta(c,L)\) obeys
\[
n(L)\eta(c,L)\to0
\]
on the intended observation horizon.
\item \textbf{Domain and interpretation.} The theorem states whether it is an
operational record theorem only or also claims an ontic one-history law.
\end{enumerate}

When these four rows are proved from the same selected source, the previous
informal slogan can be replaced by a valid implication:
\[
\begin{gathered}
\text{Born-compatible source law}\\
+\text{decoherence}\\
+\text{concentration}\\
+\text{persistence}
\end{gathered}
\Longrightarrow
\text{controlled classical record theory}.
\]

\section{Consequences and non-consequences}

The corrected separation yields several firm conclusions.

\begin{itemize}
\item Measurement can remain an ordinary physical interaction; no special
conscious observer is needed in the mathematics.
\item The current q79 binary recorder is a real advance over a purely
conditional Gleason reconstruction.
\item Classical behavior is quantitative: it comes with the three errors
\(\delta\), \(\varepsilon\), and \(n\eta\) for decoherence, concentration, and
persistence.
\item Exact decoherence does not imply a deterministic outcome.
\item A sharply concentrated law does not explain the general Born rule.
\item Operational probabilities do not, by themselves, select one ontic
history.
\item Projection can define record equivalence classes, but it does not create
their measure.
\end{itemize}

The phrase ``same problem'' should therefore be retired. The two problems are
adjacent stages of one physical workflow and can be tested on one common
instrument, but their proof obligations remain distinct.

\section{Conclusion}

Born probabilities and the classical limit share an outcome algebra, not a
single theorem. A Born source theorem fixes the weights of records. Decoherence
controls interference. Concentration makes one record approximately
deterministic, and persistence keeps it stable. One-history actualization, if
required, is another question.

MTT now closes more of this chain than version 1 reported correctly. The
canonical q79 binary one-anchor Fock recorder has an exact stopped-output law
and exact second-moment capture descent on its declared domain, without a fitted
probability. The general apparatus family remains open. The classical
concentration and persistence theorems proved here then state exactly what must
be shown, after a valid record law exists, to obtain a controlled classical
regime.

This separation is not a retreat from unification. It is the structure needed
for a rigorous one: one source may eventually discharge several adjacent proof
obligations, but none is counted as solved merely because the same record labels
appear in all of them.

% BEGIN MTT MANAGED COMPUTATIONAL EVIDENCE
\section*{Computational Evidence and Reproducibility}
The concentration and persistence bounds are proved in this paper, while the general Born-source theorem remains open. The strict-upgrade ledger is included only as a current source-level boundary; it does not provide the missing outcome weights or select one realized history.

The referenced rows are frozen to the curated results repository state identified below. Hashes are grouped in eight-character blocks for line breaking.
\begin{quote}\small
\textbf{Repository:} \url{https://github.com/PeterNero/mtt-results-repro}\\
\textbf{Commit:} \texttt{31247ebb 5c22f3fb b5443024 365433c6 ee0bff4a}\\
\textbf{Manifest:} \path{release/result_manifest.json}\\
\textbf{Manifest SHA-256:}\\
\texttt{fb399689 60b00584 631dbf53 1a708e18}\\
\texttt{ef928d6b 6d935119 c185d7f6 32b1e7cd}
\end{quote}

Tier labels are quoted verbatim from that manifest. A row used directly supports only the specific computational statement identified above; a corpus-state cross-check does not prove this paper's local theorems; and an open row is evidence of an unresolved obligation, never of closure.

\par\begingroup\dimen0=\pagegoal\advance\dimen0 by -\pagetotal\ifdim\dimen0<7\baselineskip\newpage\fi\endgroup
\paragraph{Open boundary (not evidence of closure).}
\begin{itemize}
\setlength{\itemsep}{0.25em}
\item \path{A05/strict_upgrade_ledger} (\emph{open}).\par Current 2/9 strict no-knob upgrade ledger.
\end{itemize}

No imported row changes theorem ownership or promotes a neighboring claim: all local statements retain their stated hypotheses, domains, and limitations.
% END MTT MANAGED COMPUTATIONAL EVIDENCE

\begin{thebibliography}{99}

\bibitem{Gleason1957}
A.~M. Gleason,
\emph{Measures on the Closed Subspaces of a Hilbert Space},
Journal of Mathematics and Mechanics \textbf{6} (1957) 885--893,
doi:10.1512/iumj.1957.6.56050.

\bibitem{CavesEtAl2004}
C.~M. Caves, C.~A. Fuchs, K.~Manne, and J.~M. Renes,
\emph{Gleason-Type Derivations of the Quantum Probability Rule for
Generalized Measurements},
Foundations of Physics \textbf{34} (2004) 193--209,
doi:10.1023/B:FOOP.0000019581.00318.a5,
arXiv:quant-ph/0306179.

\bibitem{Zurek2003}
W.~H. Zurek,
\emph{Decoherence and the Transition from Quantum to Classical---Revisited},
arXiv:quant-ph/0306072.

\end{thebibliography}

\end{document}
